\section{Background}
\label{sec:background}

\subsection{Command in a Power Grid}

A power company controls its substations from a control center. An operator at a console
sends orders through the SCADA system to devices in the substations. A typical order opens
or closes a breaker. One operator controls many substations. In the 2015 attack, the
attackers ``operated numerous sites under the control of the
dispatcher''~\cite[p.~14]{eisac2016ukraine}.

The same order can be correct or dangerous. An operator opens a breaker to isolate a fault
or to let personnel do work on equipment. An attacker opens the same breaker to remove
power from customers. The device gets the same order in the two conditions. An analysis
with a model shows the possible effect of such orders on many substations~\cite{ten2018impact}.

Substations also have local control. A person at the substation can operate the equipment
without the network. After the 2015 attack, the companies had ``to move to manual
operations''~\cite[p.~v]{eisac2016ukraine}. \pending{Not done at this time: a power
engineer must examine the text about local control in this paper.}

\subsection{The 2015 Attack and the Model}

Three facts of the attack are in the model of this paper. First, the attackers had the
authority of the operators. They were in the networks of the companies ``for six months or
more''~\cite[p.~3]{eisac2016ukraine}, and they used the consoles of the operators. Second,
the attackers controlled the channels. The report gives steps to ``take control of the
operator workstations and thereby lock the operators out of their
systems''~\cite[p.~8]{eisac2016ukraine}. It also records ``thousands of calls'' to the
call center of a company~\cite[p.~2]{eisac2016ukraine}.

Third, the attackers were not at the substations. The companies moved to manual operations.
An analysis of two subsequent software attacks on power grids is
in~\cite{salazar2024industroyers}.

\Cref{tab:names} gives the names that this paper uses and their equivalents in the grid
example. The model is not only for a grid. It applies to each system in which one person at
a time controls devices through a network.

\begin{table}[t]
\caption{Names in this paper and their equivalents in the grid example. A power engineer
did not examine the right column at this time.}
\label{tab:names}
\centering
\small
\begin{tabularx}{\columnwidth}{@{}l>{\raggedright\arraybackslash}X@{}}
\toprule
Name & Equivalent in the grid example \\
\midrule
Commander & The operator who has authority for an area of the grid \\
Officer & A person who can be the operator for that area, or who gives approval to an operator \\
Deputy & The officer who gets command if the commander stops \\
Device & The equipment in a substation that does the orders \\
Order & A signed message, for example to open a breaker \\
Local control & A person at the substation operates the equipment \\
\bottomrule
\end{tabularx}
\end{table}
